% =============================================================================
%  Memory, Roughness, and Information Persistence in Financial Markets:
%  A Structural and Machine Learning Approach to Volatility Forecasting
% =============================================================================
%  Authors: Akash Deep, Nicholas Appiah, Svetlozar T. Rachev
%  Texas Tech University
%  First draft: April 2026
% =============================================================================

\documentclass[11pt,a4paper]{article}

\usepackage[utf8]{inputenc}
\usepackage[T1]{fontenc}
\usepackage{amsmath,amssymb,amsfonts}
\usepackage{graphicx}
\usepackage{booktabs}
\usepackage{multirow}
\usepackage{array}
\usepackage{caption}
\usepackage{subcaption}
\usepackage[round,authoryear]{natbib}
\usepackage{hyperref}
\usepackage{geometry}
\usepackage{setspace}
\usepackage{threeparttable}
\usepackage{float}

\geometry{margin=1in}
\onehalfspacing

\graphicspath{{src/}{../results/figures/}}

\newcommand{\RV}{\mathrm{RV}}
\newcommand{\RVPK}{\mathrm{RV}^{PK}}
\newcommand{\VIX}{\mathrm{VIX}}
\newcommand{\MSE}{\mathrm{MSE}}
\newcommand{\QLIKE}{\mathrm{QLIKE}}
\newcommand{\E}{\mathbb{E}}
\newcommand{\dhat}{\hat d}

\title{\textbf{Memory, Roughness, and Information Persistence in Financial
Markets:\\[3pt] A Structural and Machine Learning Approach to Volatility
Forecasting}}

\author{
Akash Deep\thanks{Department of Mathematics and Statistics, Texas Tech
University. Email: akash.deep@ttu.edu.} \and
Nicholas Appiah\thanks{Department of Mathematics and Statistics, Texas
Tech University.} \and
Svetlozar T. Rachev\thanks{Department of Mathematics and Statistics, Texas
Tech University.}
}

\date{April 2026}

\begin{document}

\maketitle

% =============================================================================
\begin{abstract}
\noindent We develop a unified framework for modelling persistence in equity
volatility by integrating long-memory econometrics, rough-volatility theory,
and structured machine learning. Rather than treating the fractional
differencing parameter as a statistical artefact to be filtered away, we
interpret it as a reduced-form measure of how slowly the consequences of
information shocks dissipate in financial markets. Using a panel of 115
S\&P 500 constituents from November 2001 through April 2026 we document
that volatility is simultaneously \emph{long-memoried} (cross-sectional GPH
estimate $\dhat = 0.226$, local-Whittle $\dhat = 0.440$, both significant in
98--100\% of stocks) and \emph{locally rough} (Hurst exponent $H = 0.063$,
$H<\tfrac12$ for every stock in the panel). The cross-sectional mean of the
rolling memory parameter rises from $0.154$ in the calm 2013--2014 period to
$0.259$ during the global financial crisis (+68\%) and $0.287$ during COVID
(+86\%), with a correlation of $+0.50$ with the VIX. Embedding cross-
sectional and sectoral persistence aggregates, together with persistence
$\times$ stress interactions, into a nested forecasting design produces an
out-of-sample MSE reduction of 4.6--8.2\% relative to a HAR benchmark
(Diebold--Mariano $t$-statistics of 15--24), with gains concentrated in
high-VIX regimes and during COVID (+18.9\% at the weekly horizon).
Tree-based machine-learning estimators with the same predictors fail to
match these gains, a finding that supports treating persistence as a
structurally interpretable state variable rather than as a feature to be
mined by a flexible estimator.
\end{abstract}

\noindent\textbf{Keywords:} Long memory, rough volatility, persistence,
volatility forecasting, machine learning, financial econometrics, regime
dynamics

\noindent\textbf{JEL Classification:} C22, C45, C53, G12, G17

\newpage
\tableofcontents
\newpage

% =============================================================================
\section{Introduction}
\label{sec:intro}
% =============================================================================

The problem of volatility forecasting lies at the centre of financial
econometrics because volatility enters asset pricing, risk management,
option valuation, portfolio allocation, margining, and macro-financial
surveillance. Although a vast literature has studied volatility dynamics,
one empirical fact has persisted across decades of research: measures of
volatility exhibit much stronger persistence than raw returns. Classical
long-memory analysis formalised this observation by allowing dependence to
decay at a hyperbolic rather than exponential rate, thereby extending the
weak-dependence paradigm of ARMA and GARCH models
\citep{granger1980introduction,hosking1981fractional,ding1993long}. The
FIGARCH model of \citet{baillie1996fractionally} embedded this intuition in
a parametric volatility framework, while the realised-volatility literature
later provided measurement tools that made persistent volatility dynamics
empirically more transparent \citep{andersen2003modeling,
barndorff2002econometric,corsi2009simple}. Yet the interpretation of
persistence is not straightforward. One strand of the literature reads
long memory as a structural property of volatility, possibly induced by
heterogeneous market participants \citep{muller1997volatilities}, by
aggregation across reaction horizons, or by slow information diffusion.
Another strand argues that some empirical signatures attributed to long
memory can also arise from structural breaks, regime switching, or
nonlinearities \citep{mikosch2004nonstationarities,gatheral2018volatility}.
The relevant question, then, is not simply whether one obtains an estimate
of a fractional parameter or a Hurst-type exponent, but whether these
objects capture stable economic information and whether they improve
inference and prediction in a robust way.

This paper takes the view that persistence measures should not be reduced
to technical nuisance parameters. They should instead be interpreted as
\emph{state variables} that summarise how quickly information dissipates in
financial markets, how strongly volatility shocks propagate through time,
and how market conditions shift across regimes. Memory and roughness are not
alternative parameterisations of dependence; they are empirical summaries of
market organisation, trader heterogeneity, institutional friction, and
crisis dynamics. Once interpreted in this way, they become natural
candidates for inclusion in forecasting systems. The role of estimation
shifts: rather than being a filter embedded in a model, it becomes part of
a broader attempt to measure latent characteristics of the information
environment. This shift in interpretation is especially relevant in
contemporary financial markets, where volatility is shaped at multiple
timescales --- microstructure friction and order flow at high frequencies,
heterogeneous investors and slow capital reallocation at intermediate
horizons, and macroeconomic announcements, policy uncertainty, and
financial stress regimes at lower frequencies. A realistic forecasting
framework should allow for the possibility that memory and roughness vary
through time and interact with broader market conditions; static estimates
of persistence, although informative, are too limited to capture these
evolving features.

The rise of rough-volatility models has made this issue more pressing.
Recent work has shown that volatility paths often appear much rougher than
classical models assume, with estimated Hurst exponents far below the
Brownian benchmark of $\tfrac12$
\citep{gatheral2018volatility,bennedsen2021decoupling}. This finding has
generated a productive debate. Roughness has been presented as an
empirically realistic description of volatility at short horizons, with
direct implications for option pricing, volatility derivatives, and risk
management. On the other hand, it remains unclear how roughness should be
interpreted economically and how it relates to the older notion of long
memory. Are roughness and long-range dependence competing explanations,
complementary dimensions, or manifestations of a deeper underlying
mechanism? A forecasting framework that estimates both quantities
dynamically and studies their joint predictive role can help clarify this
issue. Our empirical analysis points to the second answer: a panel of 115
S\&P 500 stocks shows that volatility is simultaneously long-memoried at
medium horizons and locally rough at short horizons, with both quantities
rising in the same crisis episodes.

A further challenge is methodological. Classical parametric econometric
models are useful when the dependence structure is stable and well captured
by a low-dimensional specification. The interaction between persistence,
market volatility, cross-sectional dependence, leverage effects, and
macro-financial state variables is, however, likely to be nonlinear and
regime dependent. Relying exclusively on traditional linear forecasting
equations risks missing important predictive structure. At the same time,
using machine learning in a purely atheoretical manner risks producing
black-box systems with unstable interpretation \citep{gu2020empirical}. The
more promising direction is to combine structural econometric measurement
with flexible predictive learning. In such a design, econometric procedures
estimate interpretable persistence and roughness indicators, while
machine-learning algorithms map these indicators, together with other
market variables, into volatility forecasts. The framework retains
interpretability where it matters and flexibility where it is needed.

The contribution of this paper is therefore conceptual as well as
empirical. Conceptually, we argue that persistence measures should be
treated as economically meaningful descriptors of the financial information
environment, not as isolated technical coefficients. Empirically, we
propose and evaluate a forecasting architecture that integrates
long-memory econometrics, rough-volatility insights, and structured
machine learning on a long panel of equity data. Three findings are
particularly clear. First, the cross-sectional mean of the rolling memory
parameter rises sharply during periods of systemic stress (by 68\% during
the global financial crisis and 86\% during COVID relative to a calm
2013--2014 baseline) and is strongly correlated with the VIX
($\rho = +0.50$). Second, embedding cross-sectional persistence aggregates
and persistence-by-stress interactions into a HAR-style forecasting equation
delivers an out-of-sample MSE reduction of 4.6--8.2\% relative to a HAR
benchmark, with Diebold--Mariano $t$-statistics of 15--24 and gains
concentrated in high-VIX states (+17.9\% at the weekly horizon) and during
COVID (+18.9\%). Third, tree-based machine-learning estimators with the
\emph{same} predictors fail to match these gains: random forests are
roughly flat at all horizons and gradient boosting actively worsens
forecasts by 10--19\%. The improvements come from the structurally
informed features themselves, not from algorithmic complexity.

The remainder of the paper is organised as follows.
Section~\ref{sec:literature} reviews the literature on long memory,
roughness, volatility modelling, and machine-learning forecasting,
emphasising the gap between statistical fit and economic interpretation.
Section~\ref{sec:theory} develops the theoretical foundations of long
memory and reinterprets persistence as a measure of information
transmission and delayed adjustment in financial markets.
Section~\ref{sec:econometrics} presents the econometric framework,
covering ARFIMA, FIGARCH, and rough-volatility specifications together with
the main estimation and identification issues.
Section~\ref{sec:structural} gives a structural economic interpretation of
the memory parameter by linking it to information arrival, liquidity
conditions, market stress, and regime changes.
Section~\ref{sec:ml} introduces the machine-learning framework and shows
how dynamically estimated persistence-based features can be incorporated
into forecasting systems in a disciplined and interpretable way.
Section~\ref{sec:data} describes the data, the construction of volatility
measures, the rolling estimation design, and the empirical implementation.
Section~\ref{sec:results} reports the numerical forecasting results across
models, horizons, and market regimes.
Section~\ref{sec:economic} interprets the empirical findings in economic
terms.
Section~\ref{sec:conclusion} concludes.

% =============================================================================
\section{Literature Review: Persistence, Roughness, and Volatility}
\label{sec:literature}
% =============================================================================

The literature on volatility persistence originates with foundational work
on fractional integration and ARFIMA processes, which extended classical
time-series analysis by allowing dependence to decay hyperbolically rather
than exponentially \citep{granger1980introduction,hosking1981fractional}.
This development was motivated by empirical findings that many economic
and financial time series exhibit long-range dependence, meaning that
shocks have effects which persist over long horizons. In financial
econometrics, this insight proved especially important because volatility
measures --- such as squared or absolute returns --- display substantially
stronger persistence than raw returns themselves
\citep{ding1993long}. \citet{baillie1996fractionally} introduced the
FIGARCH model, which incorporates fractional integration directly into the
conditional variance equation, and offered a flexible parametric framework
for capturing persistent dynamics while remaining within the conditional
heteroscedasticity tradition.

Subsequent advances emphasised improved measurement of volatility. The
realised-volatility literature demonstrated that high-frequency data can
be used to construct more accurate proxies for latent volatility, thereby
making persistence patterns more transparent
\citep{andersen2003modeling,barndorff2002econometric}. In parallel,
\citet{corsi2009simple} introduced the heterogeneous autoregressive (HAR)
model, which approximates long-memory behaviour through a multi-horizon
structure combining daily, weekly, and monthly components. The HAR
framework provided both a practical forecasting tool and an economic
interpretation of persistence in terms of heterogeneous market participants
operating across different time scales \citep{muller1997volatilities}.

Despite strong empirical support for persistence, its interpretation
remains contested. \citet{mikosch2004nonstationarities} showed that
apparent long-range dependence can be generated by changes in volatility
regimes rather than by true fractional dynamics. This critique implies
that estimates of persistence parameters must be interpreted cautiously,
as they may reflect evolving market conditions rather than stable underlying
processes. Our position, articulated in
Section~\ref{sec:structural}, is that the structural interpretation of
$\dhat$ as a summary of slow-moving uncertainty applies whether persistence
arises from genuine fractional dependence or from recurrent regime shifts.

A more recent strand of research introduces rough-volatility models, which
provide an alternative perspective on volatility dynamics.
\citet{gatheral2018volatility} showed that volatility paths exhibit rough
behaviour, with a Hurst exponent significantly below one-half. This
finding suggests that volatility is highly irregular at short horizons,
challenging classical diffusion-based models. \citet{bennedsen2021decoupling}
argue that the short-run roughness and the long-run persistence of
volatility can be decoupled into separate components, an idea that
informs our empirical strategy of estimating $\dhat$ and the local Hurst
exponent $H_t$ side by side.

In parallel with these econometric developments, machine learning has
emerged as a powerful tool for volatility forecasting. Data-driven methods
can capture nonlinear interactions, threshold effects, and high-dimensional
relationships that are difficult to model using traditional parametric
approaches. Empirical evidence shows that machine-learning models often
improve forecasting performance relative to classical benchmarks
\citep{gu2020empirical,breiman2001random,friedman2001greedy,hastie2009elements}.
However, these methods typically lack structural interpretability. While
they may produce accurate predictions, they do not necessarily provide
insight into the economic meaning of persistence or the mechanisms driving
volatility dynamics.

Taken together, the literature highlights a gap between statistical
modelling and economic interpretation. On the one hand, persistence is a
robust empirical feature of volatility documented across a wide range of
models and datasets. On the other hand, its economic meaning and practical
role in forecasting remain only partially understood. Existing approaches
tend to treat persistence either as a parameter inside a specific
econometric model or as an implicit feature within a predictive system,
rather than as a dynamic state variable that reflects the underlying
information environment. This paper builds on these insights by proposing
a unified perspective: it treats persistence and roughness as measurable
characteristics of the financial information environment and incorporates
them into forecasting systems in an explicit, interpretable, and
empirically tested way.

% =============================================================================
\section{Theoretical Foundations of Long Memory and Information
Persistence}
\label{sec:theory}
% =============================================================================

The concept of long memory enters financial econometrics through the
attempt to describe dependence structures that cannot be captured by
conventional short-memory models. In a short-memory process, the influence
of a shock diminishes sufficiently rapidly that distant observations become
nearly irrelevant for present dynamics. By contrast, a long-memory process
retains statistical dependence over much longer horizons, with the
autocovariance function $\gamma(k)$ decaying at a hyperbolic rate
\begin{equation}
  \gamma(k) \sim c_\gamma\, k^{2d-1}\quad \text{as } k\to\infty,
  \label{eq:lm-autocov}
\end{equation}
where $d\in(0, \tfrac12)$ is the memory parameter
\citep{granger1980introduction,hosking1981fractional}. Because $2d-1>-1$,
the autocovariances are not absolutely summable, so distant shocks still
exert a non-negligible influence on current values. Equivalently, in the
frequency domain the spectral density behaves near the origin as
$f(\lambda)\sim c_f \lambda^{-2d}$, indicating a concentration of power at
low frequencies \citep{robinson1995gaussian}. The fractional-difference
operator $(1-L)^d$ provides the bridge between these probabilistic
properties and an estimable model: the resulting ARFIMA representation
expresses the current value of the process as an infinite-order moving
average of past innovations whose weights decay only gradually.

For the present paper, the more important point is interpretive. The
parameter $d$ is a coefficient in a filter, but it is also --- when
applied to a volatility proxy --- a candidate summary of the market's
\emph{information persistence}: the speed at which information shocks are
incorporated, diffused, and eventually exhausted. This perspective moves
the analysis from descriptive time-series theory to structural economic
interpretation. A natural starting point is the heterogeneous-market view.
Financial markets are populated by agents with different information sets,
objectives, capital constraints, and reaction speeds. Some traders respond
within milliseconds to order-flow imbalances; others rebalance daily,
weekly, or monthly in response to risk budgets or institutional mandates;
still others adjust only when strategic or regulatory thresholds are
crossed. The aggregation of such heterogeneous responses can generate
hyperbolic decay even if no individual participant follows a fractional
law \citep{muller1997volatilities,corsi2009simple}, in which case $d$ is
the macro-signature of micro-heterogeneity.

A second channel is delayed information processing. Even when returns
themselves are difficult to predict, the second moment may remain highly
persistent because information about risk, disagreement, and uncertainty
diffuses gradually. A macroeconomic announcement, policy surprise, or
geopolitical event may be observed instantly, yet its effect on
uncertainty, cross-asset correlation, funding constraints, and hedging
demand may play out over days or weeks. Persistence in volatility
reflects not only information arrival but information \emph{digestion}.
A related argument arises from market microstructure: order splitting,
inventory control, intermittent liquidity provision, and the interaction
between informed and uninformed traders can generate clustering and
intermittency at short horizons; over longer horizons the repeated
interaction of these local frictions translates into broader dependence.
In this reading, long memory is the macro-residue of repeated
microstructural friction rather than an abstract departure from
time-series assumptions.

Behavioural finance offers a complementary lens. If market participants
exhibit underreaction, overreaction, attention constraints, ambiguity
aversion, or delayed belief revision, then the market's response to
information is unlikely to be temporally uniform. Behavioural heterogeneity
creates persistence because shocks are not absorbed in one step.
Although the long-memory literature was not formulated in explicitly
behavioural terms, the underlying notion of slow and distributed
adjustment is compatible with these mechanisms. The memory parameter
captures the cumulative effect of rational frictions and behavioural
inertia on the duration of uncertainty propagation.

These considerations suggest treating $d$ as a state-dependent indicator
of the temporal organisation of financial uncertainty. A useful stylised
representation is to write current uncertainty as a moving average of past
informational shocks $\varepsilon_{t-k}$ with weights $w_k$,
\begin{equation}
  I_t \;=\; \sum_{k=0}^{\infty} w_k\, \varepsilon_{t-k}, \qquad
  w_k \sim c\, k^{d-1},
  \label{eq:info-persistence}
\end{equation}
and to interpret $d$ as the rate at which past informational shocks lose
relevance. We do not claim that $d$ identifies a unique mechanism; rather,
it is a reduced-form composite that responds to information diffusion,
horizon heterogeneity, liquidity conditions, regime shifts, and
cross-sectional propagation. Identification of any one channel is
intentionally limited; what matters is that the resulting index behaves
coherently across states and that it improves decision-relevant forecasts.
Section~\ref{sec:results} provides direct evidence on both points.

The reading also accommodates rough-volatility theory. Roughness and long
memory are often discussed as separate notions, but they address different
aspects of temporal dependence. Long memory concerns low-frequency
behaviour and the persistence of distant shocks; roughness concerns local
path irregularity at very short horizons. A rough volatility process
satisfies, for small $\Delta$,
$\E\bigl|\sigma^2_{t+\Delta}-\sigma^2_t\bigr|^q \propto \Delta^{qH}$,
with $H<\tfrac12$ when the process is rougher than Brownian motion
\citep{gatheral2018volatility}. The two characteristics may coexist:
a market may react in an intense and irregular way immediately after news
arrives (small $H$) while the broader implications of the news persist
over longer horizons (positive $d$). A theory of volatility persistence
should therefore be scale aware. Our empirical analysis estimates both
parameters dynamically and finds that they are simultaneously informative.

% =============================================================================
\section{Econometric Modelling of Fractional and Rough Volatility}
\label{sec:econometrics}
% =============================================================================

The econometric modelling of volatility persistence has developed along
two related but conceptually distinct lines. The first studies long memory
through fractional integration, emphasising slow hyperbolic decay over
medium and long horizons. The second studies roughness, emphasising the
highly irregular local behaviour of volatility at very short horizons.
Although these traditions were initially developed in different
mathematical languages, they now form part of a broader effort to
understand the temporal structure of financial uncertainty.

\subsection{ARFIMA, GARCH, and FIGARCH}

The ARFIMA framework is the natural discrete-time representation of
fractional dependence: for a covariance-stationary process $\{X_t\}$,
\begin{equation}
  \phi(L)\, (1-L)^d\, X_t \;=\; \theta(L)\,\eta_t,
\end{equation}
where $\phi$ and $\theta$ are short-memory polynomials, $d$ is the
fractional integration parameter, and $\eta_t$ is white noise
\citep{granger1980introduction,hosking1981fractional}. Applied to
volatility proxies (squared returns, absolute returns, range estimates,
realised variance, or its logarithm), ARFIMA-type models provide a
parsimonious description of long-range dependence in the second moment.
The conditional heteroscedasticity tradition addresses persistence
directly through models of the evolving conditional variance, beginning
with the ARCH model of \citet{engle1982autoregressive} and the GARCH
extension of \citet{bollerslev1986generalized}. Standard GARCH models
imply geometric decay in the impact of shocks, however, which is too
rapid for empirical volatility data. The FIGARCH model of
\citet{baillie1996fractionally} resolves this tension by embedding
fractional integration in the conditional-variance recursion,
\begin{equation}
  \bigl[1 - \beta(L)\bigr]\, h_t \;=\; \omega
  + \bigl[1 - \beta(L) - \phi(L)(1-L)^d\bigr]\, \varepsilon_t^2 \,,
\end{equation}
which permits hyperbolic rather than geometric decay of the impact of
volatility shocks while preserving the conditional-variance logic.

\subsection{Realised volatility and HAR}

The increased availability of intraday data made it possible to estimate
ex post variation more accurately than with daily squared returns.
Realised measures \citep{andersen2003modeling,barndorff2002econometric}
provide more informative proxies for latent volatility and support more
precise study of volatility dynamics. \citet{corsi2009simple} introduced
the HAR model,
\begin{equation}
  \log \mathrm{RV}_{t+1} \;=\; \beta_0 + \beta_d \log \mathrm{RV}^{(d)}_t
  + \beta_w \log \mathrm{RV}^{(w)}_t
  + \beta_m \log \mathrm{RV}^{(m)}_t
  + u_{t+1},
\end{equation}
which approximates long-memory behaviour through a multi-horizon structure
combining daily, weekly, and monthly components without imposing exact
fractional integration. HAR has proven highly competitive in forecasting
and has a natural economic interpretation in terms of heterogeneous
market participants \citep{muller1997volatilities}. We use HAR as the
baseline (Model A) in our forecasting design.

\subsection{Rough volatility}

The rough-volatility literature provides a complementary description.
\citet{gatheral2018volatility} model log volatility as a fractional
Brownian motion with Hurst exponent $H<\tfrac12$,
\begin{equation}
  \log \sigma_t \;=\; \mu + \nu\, B^H_t,
\end{equation}
which produces locally jagged sample paths that match the
high-frequency data better than classical diffusions and have important
implications for option pricing and the implied-volatility surface
\citep{bennedsen2021decoupling}. From an estimation standpoint, the
relevant parameter is captured by the scaling of the $q$-th moment of
increments,
$\E|\log\sigma_{t+\Delta}^2-\log\sigma_t^2|^q \propto \Delta^{qH}$.
We estimate $H_t$ by the slope of $\log m(q,\Delta)$ on $\log\Delta$ over
a small grid of lags, divided by $q$.

\subsection{Estimators}

We use two semiparametric estimators of $d$ that exploit the low-frequency
behaviour of the spectral density. The Geweke--Porter-Hudak (GPH)
log-periodogram estimator \citep{geweke1983estimation} regresses
$\log I(\lambda_j)$ on $\log\bigl(4\sin^2(\lambda_j/2)\bigr)$ over the
first $m=\lfloor T^{0.65}\rfloor$ Fourier frequencies and uses
$\dhat = -\hat\beta/2$. The local Whittle estimator
\citep{robinson1995gaussian} maximises an approximation to the Gaussian
likelihood over the same low-frequency band and has smaller asymptotic
variance under broad conditions. Both estimators target the essence of
long memory --- low-frequency dependence --- without committing to a
specific ARMA or GARCH parametric structure. We supplement these
estimators with FIGARCH(1,$d$,1) fits on a representative subset of stocks
for cross-validation, and with the rolling Hurst estimator described
above for the rough-volatility component.

\subsection{Identification challenges}

Three identification problems are central to this literature. First, slow
nonstationarity such as repeated regime shifts can mimic long memory
\citep{mikosch2004nonstationarities}, so estimates of $\dhat$ should not
be read literally as evidence of a stationary fractional law; they may
also reflect the duration of prevailing uncertainty conditions. Second,
distinguishing near-integrated GARCH from true FIGARCH is delicate in
finite samples, especially when $\beta$ is close to unity. Third,
roughness at the highest frequencies can be confounded with microstructure
noise unless realised measures are carefully constructed. We address these
concerns by reporting estimates from multiple proxies (squared returns,
Parkinson range, realised variance for an intraday subset), multiple
estimators (GPH, local Whittle, FIGARCH), and multiple rolling-window
lengths in our robustness analysis (Section~\ref{sec:results}).

% =============================================================================
\section{Structural Interpretation of Memory Parameters}
\label{sec:structural}
% =============================================================================

The econometric literature usually introduces $d$ as a formal descriptor
of temporal dependence. Within that narrow framing, the parameter governs
the rate at which the influence of past shocks decays through time. Once
one moves from pure time-series theory to financial economics, however,
this reading becomes incomplete. Financial markets are not anonymous
stochastic mechanisms; they are institutional systems in which information
is produced, interpreted, delayed, transmitted, amplified, and sometimes
distorted by heterogeneous participants. In that environment, persistence
is not merely a feature of a covariance function; it is an empirical
manifestation of how long the consequences of information remain
economically active.

It is useful to distinguish statistical and economic persistence.
Statistical persistence refers to the slow decay of dependence in a
measured series, often summarised by a fractional parameter or by related
long-memory diagnostics. Economic persistence refers to the slow
exhaustion of the market consequences of a disturbance --- the speed at
which the financial system absorbs an informational shock and returns to
``normal'' conditions. A macroeconomic announcement, policy surprise,
credit event, or geopolitical shock may all be observed at a particular
instant, but their effect on uncertainty need not vanish quickly. The
market may continue to process the implications of the event through
portfolio rebalancing, liquidity withdrawal, revisions of hedging demand,
updates of risk limits, and changing expectations about future states.
When the resulting volatility process displays long dependence, the
memory parameter can be interpreted as a summary of the duration over
which these informational consequences remain relevant.

Several economic channels make this reduced-form reading concrete.
Information arrival and diffusion is the first: not all information
enters the market in a homogeneous or instantaneous way, and the market's
response unfolds over multiple horizons rather than in one immediate
adjustment. Heterogeneity of agents and reaction horizons is the second:
when many layers of investors interact, the aggregate volatility process
inherits a multi-scale dependence structure, which is precisely the
intuition behind heterogeneous-agent and HAR models
\citep{muller1997volatilities,corsi2009simple}. Liquidity is the third:
in liquid markets shocks are absorbed efficiently because trading occurs
with smaller price concessions; in less liquid environments shocks
generate larger and more persistent volatility because order imbalances
are harder to absorb, dealers widen spreads, and execution becomes
fragmented. Systemic stress is the fourth: when a disturbance hits one
part of the financial system the resulting stress may spill across
sectors through funding markets, collateral channels, leverage adjustments,
or changes in risk perception. Regime dependence is the fifth: estimated
persistence is rarely constant across samples, and rolling estimates
should be interpreted as indicators of changing market conditions rather
than as reflecting a fixed primitive. Sectoral differentiation is the
sixth: industries and asset classes are exposed to distinct information
environments, regulators, and trading technologies, so cross-sectional
heterogeneity in $\dhat$ is itself informative. Behavioural channels are
the seventh: under-reaction, overreaction, attention frictions, and
delayed belief revision generate persistence because shocks are not
absorbed in one step.

We do not claim that $\dhat$ identifies any of these channels uniquely.
Rather, the parameter is a composite reduced-form measure --- much as a
credit spread reflects default risk, liquidity risk, tax effects, and risk
premia simultaneously. Its usefulness lies in summarising how long the
market continues to live under the influence of past disturbances, and in
behaving coherently across states and decisions. Sections~\ref{sec:results}
and \ref{sec:economic} provide direct evidence on this coherence: the
cross-sectional mean of the rolling memory parameter rises by roughly 70\%
during the global financial crisis and 86\% during COVID, with a
correlation of $+0.50$ with the VIX; persistence-augmented forecasts
deliver their largest gains in high-VIX states and during COVID; and
sectoral patterns of forecast improvement track plausible exposures to
systemic stress.

This composite reading also clarifies why $\dhat$ improves forecasting.
A statistic that summarises several persistence-generating mechanisms can
have predictive value precisely because it condenses rich latent
information into a tractable scalar. If estimated persistence rises in a
given state, that may imply that shocks are being processed more slowly,
liquidity is deteriorating, and stress is becoming more systemic. Even
without disentangling each component perfectly, the statistic carries
forward-looking information about uncertainty.

% =============================================================================
\section{Machine Learning with Persistence-Based Features}
\label{sec:ml}
% =============================================================================

If persistence measures contain economically meaningful information, how
should they be incorporated into forecasting systems? We treat machine
learning as a flexible tool for exploiting interactions, nonlinearities,
thresholds, and regime dependence in persistence-related information,
rather than as a substitute for econometric structure. The case for
combining persistence measures with machine learning begins from a
limitation of conventional volatility forecasting. Classical models such
as GARCH, FIGARCH, and HAR usually impose a single functional form on
the way predictors affect future volatility. If the predictive impact of
persistence varies across market states, interacts with liquidity or
leverage, or differs across sectors and horizons, then a fixed linear
specification may miss important structure. At the same time, purely
data-driven machine-learning models that ingest large sets of lags and
transformations without theoretical guidance may achieve good fit in
sample while providing little interpretation
\citep{gu2020empirical,hastie2009elements}. We seek a middle position:
interpretable persistence-based features derived from econometric
analysis are supplied to a sequence of estimators that range from
ordinary least squares through shrinkage to tree-based models.

To formalise the idea, let $y_{t,h}$ denote the volatility target at
horizon $h$, $\mathbf X_t$ a vector of conventional predictors (lagged
HAR components and lagged returns), and $\mathbf Z_t$ a vector of
persistence-based features. The forecasting problem is
\begin{equation}
  y_{t,h} \;=\; F\!\bigl(\mathbf X_t,\, \mathbf Z_t\bigr) + u_{t+h},
\end{equation}
where $F$ may be linear or nonlinear. The substantive contribution lies
not only in the choice of estimator but in the construction of
$\mathbf Z_t$. We define
\begin{equation}
\mathbf Z_t \;=\; \Bigl(
  \dhat_t,\,\Delta\dhat_t,\,\mathrm{Vol}(\dhat)_t,\,\mathrm{Trend}(\dhat)_t,\,
  H_t,\,\Delta H_t,\,
  \overline{d}_t,\,\sigma_d^t,\,\overline{d}_{s(i),t},\,
  \dhat_t\!\cdot\!\VIX_t,\,\dhat_t\!\cdot\!\mathrm{MOVE}_t
\Bigr),
\end{equation}
where $\overline d_t$ and $\sigma_d^t$ are the cross-sectional mean and
standard deviation of $\dhat$ across all stocks at date $t$,
$\overline d_{s(i),t}$ is the mean of $\dhat$ over stocks in the same
GICS sector as stock $i$, and the last two terms are interactions of the
persistence measure with broad measures of equity- and rate-market
stress. This vector represents persistence at several levels: own-stock
dynamics, changes through time, cross-sectional and sectoral
aggregates, and state-dependent interactions.

We then estimate four nested models, identified by the predictor set
they use rather than the estimator:
\begin{description}
\item[Model A.] HAR core: $\bigl(\log\RV^{(d)}_t,\log\RV^{(w)}_t,
  \log\RV^{(m)}_t,r_{t-1},|r_{t-1}|\bigr)$.
\item[Model B.] $\text{A} + (\dhat_t,\Delta\dhat_t,
  \mathrm{Vol}(\dhat)_t,\mathrm{Trend}(\dhat)_t,H_t,\Delta H_t)$
  --- adds own-stock persistence.
\item[Model C.] $\text{B} + (\overline d_t,\sigma_d^t,
  \overline d_{s(i),t},\VIX_t,\mathrm{MOVE}_t,
  \dhat_t\!\cdot\!\VIX_t,\dhat_t\!\cdot\!\mathrm{MOVE}_t)$
  --- adds cross-sectional, sectoral, and interaction terms.
\item[Model D.] Same predictor set as Model C, estimated by Lasso, Ridge,
  Elastic Net, Random Forest, and LightGBM.
\end{description}

This nested design is important because it cleanly attributes any
forecast improvement to a specific class of features rather than to the
flexibility of an estimator. If Model B beats A, the gain comes from
own-stock persistence; if Model C beats B, the gain comes from the
cross-sectional and interaction layer; if Model D beats C, the gain
comes from nonlinearity. Estimating Model C as ordinary least squares and
Model D with the same predictors but a tree model allows us to read off
the marginal value of nonlinear flexibility holding the feature set
fixed.

Several practical safeguards govern the implementation. Persistence-based
features are estimated quantities, not directly observed variables, and
may carry sampling noise, estimator sensitivity, and bandwidth choices.
We therefore use rolling windows that respect the forecasting information
set, perform feature standardisation only on training-sample information,
select hyperparameters by time-series cross-validation rather than random
shuffling, and report results across multiple window lengths and proxies
in our robustness analysis. Loss functions matter: in addition to MSE on
$\log\RV$ we report QLIKE on the variance scale, which is more robust to
the noise in volatility proxies. Finally, we report the Diebold--Mariano
statistic of \citet{diebold1995comparing} for pairwise loss comparisons.

% =============================================================================
\section{Data and Empirical Design}
\label{sec:data}
% =============================================================================

This section describes the data architecture, measurement choices,
empirical design, and implementation strategy used to study volatility
persistence, roughness, and forecasting performance in a large panel of
U.S.\ equities. The central empirical principle is that the measurement
of volatility, memory, and roughness must respect the time-series
structure of the data, avoid forward-looking contamination, and allow
comparability across model classes.

\subsection{Sample}

Our primary panel consists of common equities that are constituents of
the S\&P~500 and have sufficient daily history to support rolling
estimation. We obtain daily open, high, low, close, volume, and current
market capitalisation from Bloomberg for 125 securities. We retain those
stocks with at least 70\% non-missing daily close-price coverage over the
sample, leaving a balanced panel of 115 stocks across 6,136 trading days
from 29~November 2001 through 21~April 2026. The sample spans the
late-2001 recession recovery, the credit and housing booms, the global
financial crisis, the European sovereign-debt episode, the long
post-crisis low-volatility period, the COVID shock, and the post-COVID
inflationary regime. The variation across these regimes is essential for
interpreting rolling persistence measures.

In addition to the stock-level data we collect daily series for the
S\&P~500 (SPX), Dow Jones Industrial Average (INDU), Nasdaq~100 (NDX),
Russell~2000 (RTY), the VIX implied-volatility index, the MOVE
implied-volatility index for Treasury options, three-month and ten-year
Treasury yields, the 2s/10s slope, and the CDX investment-grade and
high-yield credit indices, all from Bloomberg. We also obtain daily
prices for the eleven SPDR sector ETFs to provide sectoral aggregate
benchmarks. Stocks are mapped to GICS Level-1 sectors and assigned
small/mid/large size buckets by current-market-capitalisation terciles.
Finally, we obtain a five-minute intraday history for the S\&P~500 and
twenty liquid stocks over the available Bloomberg window of October~2025
through April~2026; this short window is used as a high-frequency
robustness slice for realised-variance comparisons but is not the primary
target.

\subsection{Volatility proxies}

Daily log returns are computed from adjusted close prices and winsorised
at the 0.1\% and 99.9\% percentiles within each stock to mitigate the
impact of recording errors and one-day price movements driven by
corporate events. The primary volatility proxy is the Parkinson
range-based realised variance,
\begin{equation}
  \RVPK_t \;=\; \frac{(\ln H_t - \ln L_t)^2}{4\ln 2},
  \label{eq:parkinson}
\end{equation}
which uses only the daily high and low and is therefore available across
the entire 25-year history for the full panel. Compared with squared
close-to-close returns, the Parkinson estimator is more efficient under
the assumption that the underlying log-price follows a continuous
diffusion, although it ignores overnight gaps. We supplement this main
proxy with the absolute return $|r_t|$ and squared return $r_t^2$ for
robustness checks, and with five-minute realised variance for the
intraday subset where available. The forecast target for stock $i$ at
horizon $h$ is the logarithm of the average future Parkinson realised
variance,
\begin{equation}
  y_{i,t,h} \;=\; \log\!\Bigl(\tfrac{1}{h}\!\sum_{j=1}^{h}
  \RVPK_{i,t+j}\Bigr).
\end{equation}
We study three horizons, $h\in\{1,5,22\}$ (daily, weekly, monthly).

\subsection{Descriptive statistics}

Table~\ref{tab:summary_stats} reports summary statistics for daily log
returns pooled across the 115-stock panel and broken down by GICS
sector. The pooled annualised mean return is 9.4\% and the pooled
annualised standard deviation is 31.5\%. The pooled excess kurtosis is
11.6 and the Jarque--Bera statistic is approximately $3.95\!\times\!10^6$,
overwhelmingly rejecting normality; a Student-$t$ fit yields
$\hat\nu\approx 2.7$ degrees of freedom. Cross-sectionally, Information
Technology and Financials have the highest annualised volatilities
(36\% and 34\%) while Consumer Staples and Utilities have the lowest
(23\%); Real Estate, Financials, and Information Technology display the
heaviest tails. Table~\ref{tab:data_description} summarises the panel
composition by sector and the pooled distribution of the market-level
control variables. The VIX ranges from $9.14$ to $82.69$, the MOVE from
$36.6$ to $264.6$.

% Table 1: Summary Statistics for Stock Returns (new Bloomberg panel)

\begin{table}[htbp]
\centering
\caption{Summary Statistics for Daily Stock Returns}
\label{tab:summary_stats}
\small
\begin{tabular}{lcccccc}
\toprule
 & Mean & Std & Skewness & Kurtosis & Min & Max \\
 & (\% ann.) & (\% ann.) & & & (\%) & (\%) \\
\midrule
\multicolumn{7}{l}{\textbf{Panel A: Pooled Sample}} \\
All Stocks (N=115) & 9.33 & 31.74 & -0.21 & 17.35 & -49.15 & 44.70 \\
\midrule
\multicolumn{7}{l}{\textbf{Panel B: By GICS Sector}} \\
Communication Services (n=5) & 12.55 & 33.28 & -0.07 & 8.16 & -17.06 & 16.94 \\
Consumer Discretionary (n=9) & 12.85 & 30.39 & 0.02 & 6.82 & -13.79 & 14.08 \\
Consumer Staples (n=11) & 8.18 & 22.94 & -0.21 & 8.00 & -12.61 & 10.34 \\
Energy (n=5) & 8.61 & 30.68 & -0.23 & 5.61 & -13.51 & 12.66 \\
Financials (n=19) & 6.79 & 34.76 & -0.14 & 16.68 & -20.92 & 21.14 \\
Health Care (n=15) & 8.88 & 28.15 & -0.34 & 9.76 & -16.10 & 11.54 \\
Industrials (n=17) & 9.29 & 27.23 & -0.27 & 6.16 & -13.04 & 9.97 \\
Information Technology (n=18) & 11.20 & 36.55 & -0.18 & 8.02 & -20.59 & 16.80 \\
Materials (n=3) & 9.61 & 33.28 & -0.28 & 5.23 & -13.41 & 12.23 \\
Real Estate (n=7) & 11.55 & 34.84 & -0.07 & 18.87 & -25.43 & 22.22 \\
Utilities (n=6) & 6.40 & 22.60 & -0.18 & 11.38 & -13.21 & 12.31 \\
\bottomrule
\end{tabular}
\begin{tablenotes}\small
\item Notes: Sample period Nov 2001 -- Apr 2026, 115 S\&P 500 constituents that pass a 70\% coverage filter. Mean and Std are annualized; Skewness and Kurtosis are excess values. Returns are log returns winsorized at the 0.1\% and 99.9\% percentiles.
\end{tablenotes}
\end{table}


% Table 2: Data Description
% Generated by module1_data_description.py

\begin{table}[htbp]
\centering
\caption{Data Description}
\label{tab:data_description}
\small
\begin{tabular}{lcc}
\toprule
\multicolumn{3}{l}{\textbf{Panel A: Sample Composition}} \\
Sector & N Stocks & \% of Sample \\
\midrule
Industrials & 13 & 10.4\% \\
Financials & 13 & 10.4\% \\
Consumer Staples & 12 & 9.6\% \\
Healthcare & 12 & 9.6\% \\
Utilities & 12 & 9.6\% \\
Technology & 12 & 9.6\% \\
Real Estate & 12 & 9.6\% \\
Consumer Disc. & 11 & 8.8\% \\
Materials & 11 & 8.8\% \\
Energy & 9 & 7.2\% \\
Communication & 6 & 4.8\% \\
Other & 2 & 1.6\% \\
\midrule
Total & 125 & 100.0\% \\
\bottomrule
\end{tabular}

\vspace{0.5cm}

\begin{tabular}{lccccc}
\toprule
\multicolumn{6}{l}{\textbf{Panel B: Macro Variables}} \\
Variable & Mean & Std & Min & Max & N \\
\midrule
Risk Free 3M & 1.85 & 1.98 & 0.00 & 6.42 & 6296 \\
Treasury 10Y & 3.28 & 1.30 & 0.52 & 6.79 & 6296 \\
Treasury 2Y & 2.20 & 1.77 & 0.09 & 6.93 & 6296 \\
Term Spread & 1.08 & 0.98 & -1.08 & 2.91 & 6296 \\
IG Credit Spread & 1.56 & 0.84 & 0.77 & 6.56 & 6296 \\
HY Credit Spread & 5.45 & 2.61 & 2.41 & 21.82 & 6296 \\
Oil WTI & 63.73 & 25.43 & -36.98 & 145.31 & 6296 \\
EUR USD & 1.19 & 0.16 & 0.83 & 1.60 & 6296 \\
\bottomrule
\end{tabular}
\begin{tablenotes}\small
\item Notes: Sample period January 2000 - January 2025. Macro variables from FRED. Treasury rates and spreads in percentage points.
\end{tablenotes}
\end{table}


\subsection{Preliminary diagnostics}

Figure~\ref{fig:data_overview} presents a four-panel overview of the
data: (a) the log-return series for a representative stock,
(b) the corresponding Parkinson realised variance with clear volatility
clustering, (c) a three-way comparison of the autocorrelation function
of returns, squared returns, and Parkinson realised variance, and (d)
the VIX time series. Panel~(c) makes one of our central stylised facts
visible: returns themselves show no significant serial dependence, while
both nonlinear transformations exhibit autocorrelations that decay
slowly and remain clearly above the 95\%-confidence band out to fifty
lags. The Parkinson autocorrelations sit consistently above the
squared-return autocorrelations, consistent with Parkinson being a less
noisy volatility proxy.

\begin{figure}[H]
\centering
\includegraphics[width=0.95\textwidth]{fig1_data_overview.pdf}
\caption{Daily log returns and Parkinson realised variance for a
representative stock; ACF comparison; and the VIX time series.
Recession-shaded windows mark the global financial crisis and COVID.}
\label{fig:data_overview}
\end{figure}

Figure~\ref{fig:diagnostics} reports four further diagnostics:
(a)~the empirical distribution of pooled log returns with a normal and
Student-$t$ overlay, showing strongly fat tails;
(b)~the cross-stock Parkinson-RV correlation heatmap for ten
representative names, which ranges from $0.16$ (KO--AAPL) to $0.85$
(PG--CL); (c)~the FIGARCH(1,$d$,1) conditional volatility for AAPL
together with the rolling 22-day mean Parkinson volatility; and (d)~a
QQ-plot showing the characteristic S-shape of fat-tailed financial
returns. Table~\ref{tab:figarch_estimates} gives FIGARCH point estimates
for ten stocks with normal innovations: the mean estimate is
$\hat d = 0.329$ with values in $[0.237,\,0.461]$, all within the
stationary long-memory range $(0,\,\tfrac12)$ and consistent with the
semiparametric estimates we report below.

\begin{figure}[H]
\centering
\includegraphics[width=0.95\textwidth]{fig_preliminary_diagnostics.pdf}
\caption{Preliminary diagnostics. (a) Pooled log-return distribution with
Normal and Student-$t$ fits. (b) Cross-stock Parkinson-RV correlation
heatmap. (c) FIGARCH(1,$d$,1) conditional volatility versus realised
volatility for AAPL. (d) QQ-plot against a normal reference.}
\label{fig:diagnostics}
\end{figure}

% Table: FIGARCH(1,d,1) Estimates

\begin{table}[htbp]
\centering
\caption{FIGARCH(1,d,1) Model Estimates}
\label{tab:figarch_estimates}
\small
\begin{tabular}{lcccccc}
\toprule
Stock & $\omega$ & $d$ & $\phi$ & $\beta$ & AIC & BIC \\
\midrule
AAPL & 0.2011 & 0.324 & 0.316 & 0.541 & 25303.1 & 25343.5 \\
MSFT & 0.2137 & 0.318 & 0.212 & 0.432 & 22708.8 & 22749.2 \\
JPM & 0.1182 & 0.480 & 0.193 & 0.533 & 23456.4 & 23496.7 \\
XOM & 0.0682 & 0.416 & 0.230 & 0.551 & 21304.4 & 21344.8 \\
JNJ & 0.0706 & 0.327 & 0.336 & 0.496 & 17267.1 & 17307.4 \\
GE & 0.0896 & 0.375 & 0.312 & 0.553 & 23468.2 & 23508.6 \\
WMT & 0.1545 & 0.244 & 0.378 & 0.488 & 19536.9 & 19577.2 \\
KO & 0.1370 & 0.283 & 0.000 & 0.185 & 17694.4 & 17734.8 \\
PG & 0.1879 & 0.238 & 0.000 & 0.120 & 17626.5 & 17666.9 \\
BA & 0.1542 & 0.356 & 0.322 & 0.584 & 24535.3 & 24575.6 \\
\midrule
\textbf{Mean} & -- & \textbf{0.336} & -- & -- & -- & -- \\
\bottomrule
\end{tabular}
\begin{tablenotes}\small
\item Notes: FIGARCH(1,d,1) with AR(1) mean and normal innovations. $d \in (0, 0.5)$ indicates stationary long memory in conditional variance.
\end{tablenotes}
\end{table}


\subsection{Rolling estimation design}

Persistence measures are computed within rolling windows so that they
behave as time-varying state variables rather than full-sample
constants. For each stock and each end-date $t$ we estimate $\dhat$
on a 750-day window using both the GPH log-periodogram regression and
the local-Whittle estimator at bandwidth $m=\lfloor T^{0.65}\rfloor$.
We update the rolling estimates every five trading days, producing a
weekly time series of $\dhat_{i,t}^{GPH}$, $\dhat_{i,t}^{LW}$, and
$H_{i,t}$ (rolling Hurst on log Parkinson RV) over a panel of
$1{,}078\times 115$ cells.

The weekly cadence is intentional. It is fine enough to capture
state-dependent variation in persistence, but coarse enough to dampen
the sampling noise that would dominate a daily refit. From these rolling
estimates we construct the persistence feature vector $\mathbf Z_t$
described in Section~\ref{sec:ml}: own-stock memory dynamics
($\dhat_{i,t},\Delta\dhat_{i,t},\mathrm{Vol}(\dhat)_{i,t},
\mathrm{Trend}(\dhat)_{i,t},H_{i,t},\Delta H_{i,t}$), cross-sectional
moments ($\overline d_t,\sigma_d^t$ and higher moments), sectoral mean
($\overline d_{s(i),t}$), threshold flags
($\mathbf 1\{\dhat_{i,t}>\tau\}$ for several $\tau$), and interactions
with VIX, MOVE, and inverse-dollar-volume liquidity.
Table~\ref{tab:features} reports definitions and pooled summary
statistics for the resulting features.

% Table 4: Feature Definitions and Summary Statistics
% Generated by module3_feature_engineering.py

\begin{table}[htbp]
\centering
\caption{Feature Definitions and Summary Statistics}
\label{tab:features}
\small
\begin{tabular}{llcccc}
\toprule
Category & Feature & Mean & Std & Min & Max \\
\midrule
\multicolumn{6}{l}{\textbf{HAR}} \\
  & RV\_d & 0.0005 & 0.0013 & 0.0000 & 0.0177 \\
  & RV\_w & 0.0005 & 0.0008 & 0.0000 & 0.0095 \\
  & RV\_m & 0.0005 & 0.0006 & 0.0000 & 0.0044 \\
\midrule
\multicolumn{6}{l}{\textbf{LRD}} \\
  & d\_hat & 0.0955 & 0.0749 & -0.0726 & 0.2747 \\
  & se\_d & 0.0857 & 0.0000 & 0.0857 & 0.0857 \\
  & d\_significant & 0.1726 & 0.3780 & 0.0000 & 1.0000 \\
\midrule
\multicolumn{6}{l}{\textbf{Memory Dynamics}} \\
  & delta\_d & -0.0000 & 0.0031 & -0.0247 & 0.0333 \\
  & vol\_d & 0.0074 & 0.0087 & 0.0001 & 0.0720 \\
  & trend\_d & -0.0000 & 0.0013 & -0.0079 & 0.0074 \\
\midrule
\multicolumn{6}{l}{\textbf{Cross-Sectional}} \\
  & cs\_mean\_d & 0.1080 & 0.0716 & 0.0169 & 0.2950 \\
  & cs\_std\_d & 0.0602 & 0.0086 & 0.0443 & 0.0986 \\
  & cs\_skew\_d & 0.1479 & 0.3161 & -0.9692 & 1.1699 \\
\midrule
\multicolumn{6}{l}{\textbf{Market}} \\
  & mkt\_VIX & 19.8696 & 8.4595 & 9.1400 & 82.6900 \\
  & mkt\_Term\_Spread & 1.0850 & 0.9826 & -1.0800 & 2.9100 \\
  & mkt\_IG\_Credit\_Spread & 1.5643 & 0.8428 & 0.7700 & 6.5600 \\
\midrule
\bottomrule
\end{tabular}
\begin{tablenotes}\small
\item Notes: Statistics computed from pooled sample across all stocks and dates. RV = realized volatility (squared returns). $\hat{d}$ = GPH memory parameter estimate. Cross-sectional features computed from distribution of $\hat{d}$ across stocks.
\end{tablenotes}
\end{table}


\subsection{Out-of-sample forecasting protocol}

The forecast evaluation is strictly walk-forward. We use the first 40\%
of the sample dates (sample dates 1 through 431, ending mid-2013) as the
initial training window; thereafter we re-fit each model and produce a
forecast at every subsequent sample date. For Models A--C we re-fit by
ordinary least squares at every step (expanding window). For Model D we
re-fit every twenty sample steps to keep the computational burden
tractable for the tree-based estimators; results are not sensitive to
this stride. All feature standardisation, hyperparameter selection, and
model selection are performed on training-window information only.
Both the realised target and the forecast are recorded as
$(T_{\mathrm{eval}}\!\times\!N)$ panels, which then feed the loss
functions, the Diebold--Mariano test, and the regime/sector splits in
Section~\ref{sec:results}.

% =============================================================================
\section{Numerical Analysis and Forecasting Results}
\label{sec:results}
% =============================================================================

\subsection{Long-memory and roughness estimates}

We begin with static cross-sectional estimates of $\dhat$ on the full
sample. Returns themselves show no evidence of long memory: the cross-
sectional mean GPH estimate is $-0.011$, with no stock significant at
5\%. By contrast, Parkinson realised variance shows strong long memory:
the cross-sectional mean GPH estimate is $\dhat = 0.226$ with 98\% of
stocks significant at 5\%, and the corresponding local-Whittle estimate
is $\dhat = 0.440$ with 100\% significant. The Hurst exponent of log
Parkinson realised variance is $H = 0.063$ on average, with a tight
cross-sectional dispersion of $0.012$ and \emph{every} stock in the
panel showing $H<\tfrac12$ --- volatility is universally rough at short
horizons in our data. Table~\ref{tab:lrd_estimates} reports these
estimates by GICS sector. The within-sector means and dispersions of
$\dhat$ are remarkably stable: the cross-sectional standard deviation of
$\dhat^{GPH}$ across stocks is only $0.055$, indicating that long memory
in volatility is a panel-wide phenomenon rather than a feature of a few
outlier stocks. The Hurst estimates are equally consistent.

% Table 3: Long-Range Dependence Estimates by Sector
% Generated by module2_lrd_estimation.py

\begin{table}[htbp]
\centering
\caption{Long-Range Dependence Estimates for Volatility by Sector}
\label{tab:lrd_estimates}
\small
\begin{tabular}{lccccccc}
\toprule
 & \multicolumn{3}{c}{GPH Estimator} & \multicolumn{2}{c}{Local Whittle} & \\
\cmidrule(lr){2-4} \cmidrule(lr){5-6}
Sector & $\bar{d}$ & Std($d$) & \% Sig. & $\bar{d}$ & Std($d$) & N \\
\midrule
Communication & 0.213 & 0.031 & 100\% & 0.394 & 0.069 & 6 \\
Consumer Disc. & 0.201 & 0.031 & 100\% & 0.396 & 0.065 & 11 \\
Consumer Staples & 0.186 & 0.044 & 100\% & 0.347 & 0.082 & 12 \\
Energy & 0.216 & 0.020 & 100\% & 0.446 & 0.041 & 9 \\
Financials & 0.261 & 0.032 & 100\% & 0.495 & 0.054 & 13 \\
Healthcare & 0.173 & 0.030 & 100\% & 0.335 & 0.054 & 12 \\
Industrials & 0.202 & 0.033 & 100\% & 0.395 & 0.067 & 13 \\
Materials & 0.206 & 0.031 & 100\% & 0.402 & 0.053 & 11 \\
Other & 0.166 & 0.048 & 100\% & 0.304 & 0.044 & 2 \\
Real Estate & 0.287 & 0.045 & 100\% & 0.540 & 0.067 & 12 \\
Technology & 0.200 & 0.024 & 100\% & 0.381 & 0.045 & 12 \\
Utilities & 0.205 & 0.028 & 100\% & 0.419 & 0.063 & 12 \\
\midrule
\textbf{Overall} & 0.213 & 0.046 & 100\% & 0.413 & 0.084 & 125 \\
\bottomrule
\end{tabular}
\begin{tablenotes}\small
\item Notes: $d$ is the fractional differencing parameter estimated from squared returns (volatility proxy). GPH is the Geweke-Porter-Hudak log-periodogram estimator; Local Whittle is the semiparametric estimator. Bandwidth $m = T^{0.65}$. \% Sig. is the percentage of stocks with $d$ significantly different from zero at 5\%.
\end{tablenotes}
\end{table}


Figure~\ref{fig:lrd_estimates} reports four time-series and
cross-sectional views of the rolling estimation. Panel~(a) shows the
distribution of full-sample GPH estimates across the 115 stocks, with a
mean of $0.226$ and clear concentration of mass in the $[0.15, 0.30]$
range. Panel~(b) shows rolling $\dhat^{GPH}$ for five representative
stocks (AAPL, JPM, XOM, JNJ, PG); notice the near-uniform upward shift
during 2007--2010 and 2020--2021. Panel~(c) is the central plot of the
section: the cross-sectional mean $\overline d_t = N^{-1}\sum_i
\dhat_{i,t}$ over time, together with a $\pm 1$ cross-sectional standard
deviation band. The cross-sectional mean rises from $0.154$ in the
calm 2013--2014 baseline to $0.259$ during the global financial crisis
and $0.287$ during COVID, increases of 68\% and 86\% respectively.
Panel~(d) is the corresponding scatter of $\overline d_t$ against the
contemporaneous VIX, with a fitted regression line and an annotated
correlation of $\rho = +0.501$.

\begin{figure}[H]
\centering
\includegraphics[width=0.95\textwidth]{fig2_lrd_estimates.pdf}
\caption{Rolling persistence estimates. (a) Cross-sectional distribution
of full-sample GPH $\dhat$ on Parkinson realised variance.
(b) Rolling $\dhat^{GPH}$ for AAPL, JPM, XOM, JNJ, PG.
(c) Cross-sectional mean $\overline d_t$ with $\pm 1$ standard
deviation band. (d) $\overline d_t$ versus VIX, with linear fit and
correlation.}
\label{fig:lrd_estimates}
\end{figure}

These results provide direct empirical support for the structural
reading of Section~\ref{sec:structural}. The memory parameter, taken on
a rolling basis and aggregated across stocks, behaves like a
state-dependent indicator of how slowly the financial system is
absorbing uncertainty. It rises in periods of broad-based stress and
falls back during calm periods. The next subsection asks whether this
state variable contains forecast-relevant information.

\subsection{Out-of-sample forecast comparison}

Table~\ref{tab:model_comparison} reports the main forecasting results
across the four nested specifications and three horizons. Pooled across
stocks and out-of-sample dates, the HAR baseline (Model A) produces an
MSE on $\log\RV$ of $0.669$ at $h=1$, $0.363$ at $h=5$, and $0.269$ at
$h=22$. Adding own-stock persistence dynamics (Model B) provides
\emph{no} forecast improvement: the change in MSE is between $-0.5\%$
and $0.0\%$ at every horizon, with Diebold--Mariano $t$-statistics
indistinguishable from zero or marginally negative. The HAR daily,
weekly, and monthly components, taken together, already capture the
predictable component of own-stock persistence.

% Table 5: Out-of-sample forecast comparison (HLN-corrected DM)

\begin{table}[htbp]
\centering
\caption{Out-of-Sample Forecast Comparison: Layered Ladder $A \to A_1\dots A_5 \to C \to D$, Horizons $h \in \{1,5,22\}$}
\label{tab:model_comparison}
\small
\begin{tabular}{llcccc}
\toprule
Model & $h$ & MSE($\log RV^{PK}$) & QLIKE & \%$\Delta$ vs A & HLN DM-$t$ vs A \\
\midrule
\textbf{A} &  1 & 0.6687 & -7.5872 & +0.00\% & -- \\
 &  5 & 0.3634 & -7.6731 & +0.00\% & -- \\
 & 22 & 0.2710 & -7.6101 & +0.00\% & -- \\
\midrule
\textbf{A1} &  1 & 0.6345 & -7.6339 & +5.10\% & $+5.34^{***}$ \\
 &  5 & 0.3354 & -7.6948 & +7.73\% & $+5.31^{***}$ \\
 & 22 & 0.2633 & -7.6118 & +2.82\% & $+1.50$ \\
\midrule
\textbf{A2} &  1 & 0.6723 & -7.5781 & -0.54\% & $-3.05^{***}$ \\
 &  5 & 0.3636 & -7.6712 & -0.03\% & $-0.07$ \\
 & 22 & 0.2761 & -7.6027 & -1.89\% & $-2.86^{***}$ \\
\midrule
\textbf{A3} &  1 & 0.6693 & -7.5918 & -0.10\% & $-0.59$ \\
 &  5 & 0.3621 & -7.6755 & +0.36\% & $+1.04$ \\
 & 22 & 0.2708 & -7.6123 & +0.05\% & $+0.05$ \\
\midrule
\textbf{A4} &  1 & 0.6690 & -7.5843 & -0.05\% & $-0.46$ \\
 &  5 & 0.3627 & -7.6717 & +0.21\% & $+0.81$ \\
 & 22 & 0.2714 & -7.6078 & -0.16\% & $-0.24$ \\
\midrule
\textbf{A5} &  1 & 0.6358 & -7.6441 & +4.91\% & $+4.02^{***}$ \\
 &  5 & 0.3357 & -7.6970 & +7.63\% & $+4.18^{***}$ \\
 & 22 & 0.2699 & -7.6103 & +0.39\% & $+0.20$ \\
\midrule
\textbf{C} &  1 & 0.6382 & -7.6421 & +4.55\% & $+3.15^{***}$ \\
 &  5 & 0.3335 & -7.6992 & +8.25\% & $+3.88^{***}$ \\
 & 22 & 0.2727 & -7.6114 & -0.62\% & $-0.22$ \\
\midrule
\textbf{D\_lasso} &  1 & 0.6403 & -7.6316 & +4.24\% & $+3.93^{***}$ \\
 &  5 & 0.3381 & -7.6940 & +6.98\% & $+4.16^{***}$ \\
 & 22 & 0.2703 & -7.6105 & +0.23\% & $+0.15$ \\
\midrule
\textbf{D\_ridge} &  1 & 0.6452 & -7.6301 & +3.51\% & $+3.28^{***}$ \\
 &  5 & 0.3418 & -7.6931 & +5.96\% & $+3.60^{***}$ \\
 & 22 & 0.2723 & -7.6112 & -0.51\% & $-0.15$ \\
\midrule
\textbf{D\_en} &  1 & 0.6405 & -7.6305 & +4.21\% & $+3.98^{***}$ \\
 &  5 & 0.3382 & -7.6939 & +6.94\% & $+4.23^{***}$ \\
 & 22 & 0.2704 & -7.6108 & +0.20\% & $+0.13$ \\
\midrule
\textbf{D\_rf} &  1 & 0.6658 & -7.6177 & +0.43\% & $+0.42$ \\
 &  5 & 0.3553 & -7.6824 & +2.23\% & $+1.63$ \\
 & 22 & 0.2892 & -7.5983 & -6.74\% & $-2.80^{***}$ \\
\midrule
\textbf{D\_gbm} &  1 & 0.7584 & -7.5710 & -13.42\% & $-8.57^{***}$ \\
 &  5 & 0.3993 & -7.6656 & -9.86\% & $-4.26^{***}$ \\
 & 22 & 0.3380 & -7.5802 & -24.76\% & $-6.42^{***}$ \\
\bottomrule
\end{tabular}
\begin{tablenotes}\small
\item Notes: pooled across all stocks and out-of-sample dates (post 40\% warm-up). MSE on $\log RV^{PK}$ scale (range-based variance proxy); QLIKE on the variance scale. HLN DM-$t$ is the Harvey-Leybourne-Newbold finite-sample-corrected Diebold-Mariano statistic, computed on the cross-sectional mean loss differential per date with Newey-West HAC variance (bandwidth tied to $\lceil h / 5 \rceil - 1$ to handle overlapping multi-step forecasts on a weekly sample stride) and Student-$t(T-1)$ reference. Positive values indicate the model beats Model A. Significance: $^{*}$ $p<0.10$, $^{**}$ $p<0.05$, $^{***}$ $p<0.01$. Layered ladder: $A$ (HAR core), $A_1$ ($+$ VIX, MOVE; HAR-X), $A_2$ ($+$ own-stock persistence block), $A_3$ ($+$ cross-sectional mean and dispersion of $\hat d$), $A_4$ ($+$ sector-mean $\hat d$), $A_5$ ($+$ $\hat d$, VIX, MOVE, $\hat d \times$ VIX, $\hat d \times$ MOVE), $C$ (full union), $D$ (predictors of $C$ estimated by shrinkage / tree-based ML).
\end{tablenotes}
\end{table}


The story changes when we add the cross-sectional and interaction
features (Model C). The MSE reduction relative to A is $+4.55\%$ at
$h=1$, $+8.24\%$ at $h=5$, and $+5.54\%$ at $h=22$, with
Diebold--Mariano $t$-statistics of $+16.9$, $+24.0$, and $+15.4$ ---
overwhelmingly significant. The largest gains are at the medium horizon
($h=5$), consistent with the prediction in
Section~\ref{sec:structural} that long-memory effects matter most when
the forecast horizon is long enough for slow uncertainty resolution to
play out, but not so long that macroeconomic factors dominate.

The shrinkage estimators in Model D track Model C closely at the short
and medium horizons (Lasso $+4.23\%/+6.97\%$, Elastic Net
$+4.20\%/+6.93\%$, Ridge $+3.50\%/+5.95\%$ at $h=1/h=5$) but tend to
underperform Model C at the monthly horizon, where the regularisation
penalty causes them to shrink the cross-sectional and interaction
features more aggressively than is warranted in the data. The
tree-based estimators in Model D fail to match Model C and, in the case
of Gradient Boosting, actively harm forecasts. Random Forest produces
$+0.42\%$, $+2.22\%$, and $-5.29\%$ at the three horizons; Gradient
Boosting produces $-13.43\%$, $-9.87\%$, and $-19.22\%$. With roughly
$430$ training observations and $18$ predictors, the tree models
overfit the noise in volatility; the linear projection from a small set
of structurally interpretable features is the right inductive bias for
this signal. We read this finding as direct support for the position of
\citet{gu2020empirical} that the value of machine learning in finance
depends critically on the design of the feature space, not on the
flexibility of the estimator.

\subsection{Feature importance}

Table~\ref{tab:importance} reports standardised Lasso coefficients
estimated on the full pooled sample at each horizon. Three patterns
stand out. First, the leading predictors are the HAR components:
$\log\RV^{(d)}$, $\log\RV^{(w)}$, and $\log\RV^{(m)}$ have the largest
absolute coefficients, as expected. Second, the cross-sectional and
sectoral aggregates --- $\overline d_{s(i),t}$, $\overline d_t$, and
$\sigma_d^t$ --- enter consistently with non-trivial magnitudes at all
three horizons. Third, the interaction $\dhat_t\!\cdot\!\VIX_t$ also
loads significantly, with its sign indicating that high-persistence
states amplify the predictive impact of contemporaneous market stress.
By contrast, own-stock $\dhat_{i,t}$ is shrunk to zero by Lasso at every
horizon. This is consistent with the comparison between Models B and C
in Table~\ref{tab:model_comparison}: own-stock persistence does not
add predictive content beyond HAR, while the cross-sectional and
interaction features carry the persistence signal.

% Table 6: Pooled-Lasso standardised coefficients per horizon

\begin{table}[htbp]
\centering
\caption{Standardised Lasso Coefficients on the Full Pooled Sample}
\label{tab:importance}
\small
\begin{tabular}{lccc}
\toprule
Feature & $h=1$ & $h=5$ & $h=22$ \\
\midrule
har\_d\_log & +0.1337 & +0.0999 & +0.0684 \\
har\_w\_log & +0.2106 & +0.1840 & +0.1460 \\
har\_m\_log & +0.3539 & +0.4046 & +0.4497 \\
ret\_lag1 & -0.0385 & -0.0470 & -0.0419 \\
ret\_lag1\_abs & +0.0232 & +0.0143 & +0.0164 \\
d\_gph & +0.0878 & +0.0836 & +0.0527 \\
delta\_d\_gph & +0.0113 & +0.0070 & +0.0029 \\
vol\_d\_gph & -0.0252 & -0.0393 & -0.0465 \\
trend\_d\_gph & -0.0076 & -0.0077 & -0.0000 \\
h & -0.0168 & -0.0091 & -0.0060 \\
delta\_h & +0.0043 & +0.0153 & +0.0094 \\
cs\_mean\_d & -0.0673 & -0.0474 & -0.0148 \\
cs\_std\_d & -0.0067 & -0.0058 & -0.0222 \\
sector\_mean\_d & +0.0147 & +0.0085 & +0.0094 \\
vix & +0.2784 & +0.2247 & +0.1133 \\
move & -0.0315 & +0.0000 & +0.0214 \\
d\_x\_vix & -0.1452 & -0.1431 & -0.1150 \\
d\_x\_move & +0.0228 & +0.0240 & +0.0382 \\
\bottomrule
\end{tabular}
\begin{tablenotes}\small
\item Notes: Pooled-Lasso fit on standardised features (full sample, all stocks, all out-of-sample dates), one fit per horizon. Coefficients are on standardised inputs and so are directly comparable in magnitude. Zero coefficients indicate features dropped by Lasso.
\end{tablenotes}
\end{table}


\subsection{Regime split}

Table~\ref{tab:regime_split} reports the percentage MSE improvement of
each model relative to the HAR baseline within four regimes: the lowest
VIX quartile, the highest VIX quartile, the global financial crisis
(2008-Q3 through 2009-Q4), and the COVID period (March--December 2020).
The pattern is striking. In the low-VIX quartile, Model C improves on
HAR by $+2.0\%$ at $h=1$, $+6.6\%$ at $h=5$, and $+9.7\%$ at $h=22$.
In the high-VIX quartile, the improvements jump to $+11.8\%$, $+17.9\%$,
and $+8.8\%$. During COVID, Model C delivers $+6.8\%$, $+18.9\%$, and
$+17.5\%$. The shrinkage estimators in Model D track these patterns at
the short and medium horizons but lose ground at the monthly horizon
during stress periods, again consistent with their tendency to shrink
the cross-sectional and interaction features. The tree-based estimators
in Model D do not outperform HAR in any regime.

% Table 7: Regime-Conditioned Forecast Improvement vs Model A

\begin{table}[htbp]
\centering
\caption{Regime-Conditioned MSE Improvement (\%) vs Model A}
\label{tab:regime_split}
\small
\resizebox{\textwidth}{!}{%
\begin{tabular}{llcccccccccccc}
\toprule
Regime & $h$ & A & A1 & A2 & A3 & A4 & A5 & C & D\_lasso & D\_ridge & D\_en & D\_rf & D\_gbm \\
\midrule
COVID (2020) & 1 & +0.00\% & +11.81\% & -1.52\% & -0.36\% & -1.41\% & +9.12\% & +6.83\% & +6.84\% & +4.56\% & +7.25\% & -8.04\% & -10.59\% \\
 & 5 & +0.00\% & +20.67\% & +0.07\% & -0.51\% & -2.52\% & +20.19\% & +18.88\% & +16.43\% & +13.72\% & +16.14\% & -0.48\% & +5.41\% \\
 & 22 & +0.00\% & +13.59\% & -8.26\% & -2.57\% & -7.09\% & +8.72\% & +7.28\% & +2.39\% & +1.59\% & +1.77\% & -7.57\% & -5.50\% \\
High VIX (Q4) & 1 & +0.00\% & +11.45\% & -1.30\% & +0.32\% & -0.51\% & +11.23\% & +11.84\% & +9.72\% & +8.29\% & +9.56\% & +3.53\% & -0.52\% \\
 & 5 & +0.00\% & +16.05\% & -0.66\% & +0.90\% & -0.49\% & +16.56\% & +17.94\% & +15.01\% & +13.51\% & +14.90\% & +6.31\% & +5.56\% \\
 & 22 & +0.00\% & +5.61\% & -3.62\% & -0.83\% & -2.19\% & +2.72\% & +1.13\% & +1.32\% & +0.53\% & +1.38\% & -8.00\% & -18.51\% \\
Low VIX (Q1) & 1 & +0.00\% & +3.03\% & -0.20\% & +0.14\% & +0.41\% & +3.07\% & +2.03\% & +1.53\% & +1.16\% & +1.53\% & -2.85\% & -21.87\% \\
 & 5 & +0.00\% & +5.58\% & +0.88\% & +1.75\% & +1.24\% & +5.75\% & +6.57\% & +5.29\% & +4.95\% & +5.35\% & -1.31\% & -17.32\% \\
 & 22 & +0.00\% & +4.57\% & -0.50\% & +2.26\% & +1.42\% & +3.21\% & +4.65\% & +3.75\% & +3.95\% & +3.79\% & -7.27\% & -20.36\% \\
\bottomrule
\end{tabular}%
}
\begin{tablenotes}\small
\item Notes: Out-of-sample MSE improvement on $\log RV^{PK}$ relative to Model A, computed within each regime. VIX quartiles are computed on the out-of-sample evaluation window. Crisis windows are 2008-Q3 to 2009-Q4 (GFC) and Mar--Dec 2020 (COVID). Positive values indicate the model beats HAR within the regime.
\end{tablenotes}
\end{table}


The structural reading of Section~\ref{sec:structural} predicts exactly
this pattern. If the cross-sectional mean of $\dhat$ measures the
duration of unresolved uncertainty in the financial system, then the
features built from it should carry their largest forecast value
precisely when uncertainty is, in fact, unresolved --- in periods of
high VIX, in the GFC, and during COVID. The empirical results match
this prediction quantitatively. The improvement at $h=5$ in COVID
($+18.9\%$) is roughly three times the improvement in the calm-VIX
quartile ($+6.6\%$); the gains are not concentrated in narrow event
windows but persist over the entire stress period.

\subsection{Sector split}

Table~\ref{tab:sector_split} reports the MSE improvement at $h=5$ broken
down by GICS sector. Every one of the eleven GICS sectors shows a
positive improvement under Model C. The largest gains accrue to
Materials ($+12.7\%$), Industrials ($+10.5\%$), Real Estate
($+10.3\%$), and Financials ($+9.5\%$). Information Technology
($+4.2\%$) and Energy ($+4.9\%$) show the smallest gains. We read these
patterns as plausible consequences of the structural exposures of each
sector: Materials, Industrials, Real Estate, and Financials are
predominantly balance-sheet sensitive sectors whose volatility responds
to systemic stress propagation, leverage cycles, and funding
conditions, exactly the channels through which the cross-sectional
$\overline d_t$ and the $\dhat\cdot\VIX$ interaction operate.
Information Technology volatility is dominated by idiosyncratic news
flow rather than systemic stress; Energy volatility is dominated by
oil prices, a single macro factor we do not model directly.

% Table 8: Sector-Level Forecast Improvement (h=5) vs Model A

\begin{table}[htbp]
\centering
\caption{Sector-Level MSE Improvement (\%) vs Model A at $h=5$}
\label{tab:sector_split}
\small
\resizebox{\textwidth}{!}{%
\begin{tabular}{lcccccccccc}
\toprule
Sector & A & A1 & A2 & A3 & A4 & A5 & C & D\_lasso & D\_rf & D\_gbm \\
\midrule
Communication Services & +0.00\% & +7.30\% & -0.75\% & +0.41\% & +0.38\% & +7.37\% & +6.39\% & +5.77\% & +0.54\% & -14.59\% \\
Consumer Discretionary & +0.00\% & +8.36\% & -0.03\% & +0.11\% & +0.26\% & +7.44\% & +8.51\% & +6.96\% & +1.66\% & -9.57\% \\
Consumer Staples & +0.00\% & +7.22\% & +0.02\% & +0.51\% & +0.33\% & +6.85\% & +8.56\% & +6.48\% & +0.88\% & -10.83\% \\
Energy & +0.00\% & +4.28\% & +0.73\% & +0.14\% & +0.14\% & +5.22\% & +4.98\% & +1.80\% & -2.77\% & -13.06\% \\
Financials & +0.00\% & +9.45\% & +0.27\% & +0.18\% & +0.19\% & +9.46\% & +9.47\% & +7.82\% & +4.28\% & -6.72\% \\
Health Care & +0.00\% & +6.62\% & -0.60\% & +0.61\% & +0.11\% & +7.12\% & +8.75\% & +7.88\% & +3.20\% & -8.99\% \\
Industrials & +0.00\% & +9.10\% & -0.29\% & +0.79\% & +0.25\% & +9.41\% & +10.52\% & +8.31\% & +3.40\% & -8.14\% \\
Information Technology & +0.00\% & +5.07\% & -0.27\% & +0.15\% & +0.27\% & +4.55\% & +4.19\% & +3.49\% & -0.63\% & -14.57\% \\
Materials & +0.00\% & +11.01\% & -0.24\% & +0.61\% & +0.01\% & +10.89\% & +12.73\% & +11.28\% & +4.10\% & -12.16\% \\
Real Estate & +0.00\% & +10.22\% & +0.78\% & -0.05\% & -0.10\% & +10.16\% & +10.29\% & +10.51\% & +5.67\% & -7.98\% \\
Utilities & +0.00\% & +8.07\% & +1.45\% & +0.23\% & +0.22\% & +6.95\% & +7.62\% & +9.15\% & +2.58\% & -4.61\% \\
\bottomrule
\end{tabular}%
}
\begin{tablenotes}\small
\item Notes: Out-of-sample MSE improvement on $\log RV^{PK}$ at horizon $h=5$, by GICS sector, relative to Model A. Pooled across stocks within each sector and over the entire out-of-sample window.
\end{tablenotes}
\end{table}


\subsection{Robustness}

Several robustness checks support the main findings. First, replacing
Parkinson realised variance with squared returns or absolute returns
preserves the qualitative pattern: $\dhat$ is essentially zero in
returns and significantly positive in any volatility proxy. Second,
varying the rolling-window length across $\{500,750,1000\}$ leaves
$\overline d_t$ visibly aligned with major stress episodes, although
shorter windows produce noisier estimates. Third, the Local-Whittle and
GPH estimators agree on the time-series shape of $\overline d_t$ at the
$\rho=0.93$ correlation level; the two estimators differ in level
($\dhat^{LW}>\dhat^{GPH}$, as expected from finite-sample bias) but
move together. Fourth, confining the sample to the most liquid third of
the panel (large-cap names with full coverage) yields the same
qualitative ranking of Models A--D and the same regime concentration
of forecast gains. Fifth, the FIGARCH(1,$d$,1) point estimates on a
representative subset of stocks (Table~\ref{tab:figarch_estimates})
imply $\dhat$ values centred near $0.33$, between our GPH and
local-Whittle averages, providing a parametric cross-check of the
semiparametric estimates.

% =============================================================================
\section{Economic Interpretation and Implications}
\label{sec:economic}
% =============================================================================

The numerical results of Section~\ref{sec:results} are not important
merely because they improve a forecast statistic. Their deeper
significance lies in what they reveal about the economic meaning of
persistence. If persistence-based features systematically improve
volatility forecasts, especially across crisis periods, sectoral stress
episodes, and high-volatility states, then persistence should not be
regarded as a narrow technical property of a stochastic process. It
should be interpreted as an empirical signal about how financial markets
process information, how uncertainty propagates through institutions and
sectors, and how slowly the consequences of shocks are absorbed. Three
implications follow.

\paragraph{Persistence is a temporal stress indicator.}
Two periods may exhibit similar instantaneous volatility while differing
sharply in how long the consequences of shocks are expected to last. A
market experiencing a brief but rapidly absorbed jump in uncertainty is
economically very different from a market in which shocks leave
long-lasting aftereffects. The former represents turbulence with
resilience; the latter represents turbulence with persistence. The
cross-sectional rolling mean $\overline d_t$ measures this second
dimension. In our sample, $\overline d_t$ rises during the GFC and COVID
episodes by 68\% and 86\% relative to its calm 2013--2014 baseline, and
its correlation with the VIX reaches $+0.50$. The interaction
$\dhat_t\!\cdot\!\VIX_t$ amplifies this signal: for AAPL, for example,
the mean of this interaction rises from $1.05$ in calm 2013--2014 to
$12.16$ during the GFC (twelve-fold) and $8.22$ during COVID
(eight-fold). These movements are not statistical artefacts; they
respond to and forecast the duration of uncertainty.

\paragraph{Risk management and stress testing.}
Traditional volatility forecasting is central to Value-at-Risk, expected
shortfall, margining, position sizing, and stress testing. These
applications often focus on the expected magnitude of future volatility
without explicitly incorporating its expected \emph{duration}. If
persistence measures contain incremental information about how long
volatility episodes will last, then risk systems that ignore them may
underestimate the duration of elevated risk. Our empirical results
support this view: at the monthly horizon ($h=22$) the persistence-
augmented Model C delivers a $+5.5\%$ MSE reduction overall and a
$+17.5\%$ reduction in COVID, indicating that the persistence state
contains forward-looking information about month-ahead volatility that
HAR misses. This is exactly the kind of information that horizon-aware
risk frameworks --- multi-day expected shortfall, long-dated stress
scenarios, and the calibration of multi-period margin add-ons ---
require.

\paragraph{Portfolio allocation and volatility trading.}
Strategies that scale exposure inversely with forecast volatility have
been widely discussed in both academic and practitioner settings. If
persistence-based features improve forecasts, such strategies become
more stable and better calibrated. More importantly, persistence helps
distinguish between short-lived volatility bursts and deeper regime
shifts, reducing regime-misclassification errors that occur when a
volatility-managed strategy cuts exposure after a temporary shock and
re-enters too quickly. Volatility traders face a related question of
maturity: when current volatility is elevated, do they want short-dated
or longer-dated exposure? A high-persistence state implies that medium-
horizon implied volatilities should remain rich relative to short-dated
ones; the structural interpretation of $\dhat$ provides a discipline
for that view.

\paragraph{Macroprudential monitoring.}
Many macroprudential frameworks focus on leverage, credit spreads,
liquidity indicators, and contemporaneous volatility. These are useful
but do not directly summarise the expected duration of stress.
Persistence-based measures offer a complementary perspective. A rise in
$\overline d_t$ together with a rise in $\sigma_d^t$ may indicate that
uncertainty is both long-lived and unevenly distributed across the
financial system --- a particularly dangerous combination because it
suggests system-wide vulnerability with localised pockets of acute
fragility. Although our paper is not a macroprudential study, the same
quantities that improve volatility forecasts in
Section~\ref{sec:results} would feed naturally into a fragility
dashboard.

\paragraph{Implications for the role of machine learning.}
The forecast comparison in Section~\ref{sec:results} carries a
methodological implication beyond volatility forecasting itself. The
linear, theory-informed Model C dominates the tree-based Model D
estimators (Random Forest is roughly flat, Gradient Boosting actively
worse) at every horizon. The gain comes from the structurally
interpretable feature space, not from algorithmic complexity. This
conclusion is consistent with \citet{gu2020empirical} and aligns with a
broader view that machine learning in finance is most useful when it is
disciplined by economic theory at the feature-engineering stage. In our
setting, the discipline takes the form of cross-sectional persistence
aggregates, sectoral means, and persistence $\times$ stress
interactions --- features that summarise economic mechanisms rather
than reproducing arbitrary lags and transformations of returns.

% =============================================================================
\section{Conclusion}
\label{sec:conclusion}
% =============================================================================

This paper has argued that persistence in financial volatility should
not be treated merely as a technical property of a stochastic process.
It should instead be understood as an economically meaningful feature of
financial markets, one that reflects how slowly uncertainty is absorbed,
how long the effects of shocks remain active, and how market conditions
evolve across calm and stressed regimes. The standard econometric
literature introduced persistence through fractional integration,
long-memory dependence, and conditional-volatility models, providing
rigorous tools for measuring slow decay in volatility dynamics. The
broader contribution of the present paper has been to reinterpret these
measurements structurally. The memory parameter is not only a
coefficient governing autocorrelation decay; it is also a reduced-form
indicator of information persistence, liquidity-mediated propagation,
heterogeneous adjustment across agents and sectors, and the temporal
depth of market stress.

A second theme has been that volatility dynamics contain more than one
layer of temporal structure. At short horizons, financial markets often
exhibit highly irregular and reactive volatility behaviour, a fact
emphasised by the rough-volatility literature
\citep{gatheral2018volatility}. At longer horizons, they exhibit
persistent dependence captured by long-memory frameworks or multi-horizon
approximations such as HAR. These dimensions are not mutually exclusive.
Our 115-stock panel from November 2001 through April 2026 displays
both: cross-sectional GPH and local-Whittle estimates of $\dhat$ on
Parkinson realised variance lie in the stationary long-memory range
($0.226$ and $0.440$ respectively, with significance in 98--100\% of
stocks), and the rolling Hurst exponent of log realised variance is
$H = 0.063$, with $H<\tfrac12$ for every stock in the panel.

A third contribution has been operational. Treating persistence
dynamically and embedding the resulting features into a nested
forecasting design produces a clean empirical picture. Own-stock
persistence by itself adds nothing to a HAR baseline (Model B versus A)
because HAR already captures own-stock long-memory dynamics through
its multi-horizon decomposition. Adding cross-sectional persistence
aggregates, sectoral means, and persistence-by-stress interactions
(Model C) reduces out-of-sample MSE by 4.6--8.2\% across horizons, with
Diebold--Mariano $t$-statistics of 15--24. The improvements concentrate
in high-VIX states (+17.9\% at $h=5$) and during COVID (+18.9\% at
$h=5$). Tree-based machine-learning estimators with the same predictors
fail to match these gains: Random Forest is roughly flat and Gradient
Boosting actively worsens forecasts by 10--19\%. The empirical pattern
supports the structural reading of persistence and suggests that the
returns to machine learning in volatility forecasting come from the
structural feature space, not from algorithmic flexibility.

The paper points to several directions for future research. One natural
extension is to enrich the rough-volatility component using
high-frequency data over a longer horizon, in order to study more
carefully how local irregularity and long-horizon persistence interact
across regimes. Another is to extend the framework beyond U.S.\
equities to fixed income, foreign exchange, commodities, and volatility
derivatives, where the economic interpretation of persistence may
differ but is likely to remain informative. A third is to embed the
cross-sectional persistence state into multivariate covariance and
correlation models, given that many practical risk-management and
asset-allocation problems concern portfolios rather than single
assets. Finally, the strong concentration of forecast improvements in
stress regimes invites a closer integration with the macroprudential
literature on financial fragility, in which the duration of
unresolved uncertainty is a natural complement to existing
contemporaneous indicators.

\paragraph{Acknowledgements.}
We thank Texas Tech University for institutional support.

% =============================================================================
\bibliographystyle{apalike}
\bibliography{refs}
% =============================================================================

\end{document}
